A third check uses the local TEV decision model to assign the highest-probability U.S. occupation, construct a probability-weighted exposure conditional on a fixed 24-candidate cosine shortlist, and predict exposure directly from Spanish occupational information. Online Appendix Table~\ref{tab:v1_jev_robustness} reports the results. In the preferred within-CNO1 specification, the highest-probability and shortlist-weighted measures imply unemployment increases of approximately 1.6 and 1.7 percent during the adjustment period and 2.2 and 2.6 percent later, per 10 percentage points of the respective exposure measure. The direct score gives larger coefficients, approximately 6.0 and 14.2 percent, on a different model-predicted scale. Contract estimates are small and imprecise for all three measures. The full-sample pre-treatment joint tests reject for both outcomes under each measure, and excluding 2021 does not resolve every diagnostic. These checks therefore show a positive unemployment association under alternative constructions, while retaining measurement and identification concerns. Online Appendix Figures~\ref{fig:v1_jev_exposure_unemployed} and~\ref{fig:v1_jev_exposure_contracts} report the dynamics.
